\subsection{Sharing changes the complexity}
\label{app:proofs-circuits}

A circuit is a finite acyclic graph of operation gates and input
variables, with a designated output; a gate's value may be used more
than once. Its size counts gates and edges. Unfolding it gives a tree
term with the same value, but that tree need not have polynomial size.

\begin{theorem}[Reset circuit equality]\label{reset:thm:circuits}
For every fixed $K\geq3$, equality of reset circuits is coNP-complete,
without constants. It is already coNP-complete on the three-element
generator $Z\subset U\subset T$. The same conclusion holds for
circuits over $\{\del,H_a\}$, over $H_a$ alone, or over any nonempty
family of rooted blocking-policy operations.
\end{theorem}
\begin{proof}
A failed valuation is a polynomial certificate. For reset alone,
\cref{reset:thm:three} applies to the unfolded terms, so inputs from
$\{Z,U,T\}$ suffice; evaluate the original circuit without unfolding.

For hardness take a CNF $\varphi$ with rails $p_i,n_i$, a root $z$,
and a marker $c$. Write $\operatorname{rc}$ for the right reset comb
of \cref{app:proofs-hard}. Put
\[
 W=\operatorname{rc}(p_1,n_1,\ldots,p_n,n_n;z),\qquad
 R=\operatorname{rc}(((c\del p_i)\del n_i)_{i=1}^n;W).
\]
Enumerate the clauses as $C_1,\ldots,C_q$, reading a literal as its
rail. Set $A_0=c$ and, for $1\leq j\leq q$, define the circuit gate
\[
 A_j=\operatorname{rc}((A_{j-1}\del\ell)_{\ell\in C_j};z).
\]
All occurrences of $A_{j-1}$ refer to the same earlier circuit output.
The equation is $A_q\del R=R$. The construction has
$5n+2\sum_j|C_j|+1$ reset gates, including the final comparison's
left output, and $O(n+\sum_j|C_j|+q)$ description size.
An empty clause gives $A_j=z$.

If $z$ is atomic, $R=z$ and the equation holds. Otherwise $R$ is
ambiguous and contains $W$. For a marker $w\notin W$, the clause
recursion gives
\[
 w\in A_j\quad\Longleftrightarrow\quad
 w\in A_{j-1}\ \land\
 \exists\ell\in C_j:\operatorname{Amb}(\ell).
\]
Indeed each literal rail omits $w$, and its reset gate retains $w$
exactly when that rail is ambiguous; the right comb takes their union
with $z$. Meanwhile $w\in R$ exactly when $w\in c$ and some pair
$p_i,n_i$ is simultaneously ambiguous. Every output contribution
outside $W$ comes from $c$. A failed equation therefore gives an
ambiguous literal in every clause and never two ambiguous opposite
rails. Assign variable $i$ true iff $p_i$ is ambiguous; this satisfies
the CNF.

Conversely, use the fixed chain of \cref{reset:thm:three}. Given a
satisfying assignment, set true rails to $U$, false rails to $Z$,
$z=U$, and $c=T$. Then $W=R=U$, while every clause preserves the
extra winner of $T$, so $A_q=T$ and the equation fails. Empty
conjunctions and empty clauses are included in the construction.

Finally $x\del y=H_a(y,x,x)=R_D(y,x,x)$ for every policy $D$.
Replacing a reset gate by either of these gates uses repeated edges,
not repeated copies of its input circuit, and has constant cost.
This transfers hardness. Direct evaluation over the fixed finite
carrier gives the coNP upper bounds for all the listed signatures.
\end{proof}

Thus the guard algorithms are polynomial in tree size and cannot in
general be made polynomial in circuit size unless $\mathrm P=\mathrm{NP}$.
This is compatible with a finite equational basis: reset itself has
one. Sharing changes succinctness, not the valid identities.
\path{lean/Fragments.lean} checks the generic clause and CNF marker
recurrences as \path{circuitClause_marker} and \path{circuitCNF_marker};
syntax assembly and the complexity argument above are written proofs.
\path{checks/check_reset_circuits.py} evaluates the actual acyclic
circuits on all winner-set valuations in its finite cases and checks
the fixed-chain reduction separately.
